\section{Data leakage: the silent career killer}

\subsection{What it is}

Leakage = the model sees information during training or evaluation that it
would not have at forecast time. The backtest looks brilliant. Production
fails. Nobody knows why.

\subsection{The three leaks in load forecasting}

\begin{enumerate}
  \item \textbf{Lag leakage.} Using yesterday's full load (lag 24 h) for a
        day-ahead forecast. At the 09:00 cutoff most of yesterday is not
        measured yet. Our code raises an error if a lag crosses the cutoff.
  \item \textbf{Weather leakage.} Backtesting with \emph{actual} weather (ERA5
        reanalysis) as input. The real desk only had a weather \emph{forecast}.
        Actuals remove weather-forecast error and inflate accuracy.
        Fix: evaluate with archived forecasts at 1--2 day lead time
        (Open-Meteo Previous Runs API).
  \item \textbf{Preprocessing leakage.} Fitting scalers, feature selection, or
        early stopping on data that includes the test period. Everything is
        fitted inside the training window only, then applied forward.
\end{enumerate}

\subsection{How we prove absence, not just claim it}

Tests, not promises (\texttt{tests/test\_features.py}):

\begin{itemize}
  \item Corrupt every data point after the cutoff with garbage ($10^9$).
        Recompute features. They must be byte-identical. If a feature changes,
        it was reading the future.
  \item Ask for lag 24 explicitly. The code must raise, not comply.
\end{itemize}

The corruption test is the strong one: it catches leaks we did not think of,
because \emph{any} post-cutoff read changes the output.

\subsection{Interview line}

``I do not argue my pipeline is leak-free. I corrupt all post-cutoff data with
garbage and show the forecast does not change. That is proof, not opinion.''
